\chapter{Algebraic geometry and~analytic~geometry}
\label{XII}
% original p. 311

\begin{center}
{Mme M.\ \textsc{Raynaud}\footnote{After unpublished notes of
A\ptbl Grothendieck.}}
\end{center}
\vspace*{1cm}

Proceeding as in \cite{XII.10}, one associates with every scheme $X$
locally of finite type over the field of complex numbers $\CC$ an
analytic space~$X^\an$
whose underlying set is $X(\CC)$.

In nos.~\Ref{XII.2} and~\Ref{XII.3} of this expos\'e, we give a
``dictionary'' between the usual properties of~$X$ and of
$X^\an$ and between the properties of a morphism $f\colon X\to
Y$ and of the associated morphism $f^\an\colon X^\an \to Y^\an$.

We then show that the comparison theorems between
coherent sheaves on~$X$ and $X^\an$, established in \cite[\No
12]{XII.10} for a projective variety, are still valid when
$X$ is a proper scheme.

Finally we prove in \No \Ref{XII.5} the equivalence of the category of
finite \'etale coverings of~$X$ and of the category of
finite \'etale coverings of~$X^\an$. As a bonus for the reader, we
give a new proof of the theorem of
Grauert-Remmert \cite{XII.6}, using the resolution of
singularities~\cite{XII.8}.

\section{Analytic space associated with a scheme}
\label{XII.1}
% original p. 312

Let $X$ be a scheme locally of finite type over~$\CC$. Let $\Phi$ be
the functor from the category of analytic spaces \cite[\No 9]{XII.4} to
the category of sets which associates with an analytic space~$\othercal{X}$ the set of morphisms of spaces ringed in
$\CC$-algebras $\Hom_\CC (\othercal{X},X)$. We have the following
theorem:

\begin{theoremdefinition}
\label{XII.1.1} The functor $\Phi$ is representable by an
analytic space $X^\an$ and a morphism $\varphi\colon X^\an\to
X$. One says that $X^\an$ is the analytic space associated with~$X$.

If $|X^\an|$ is the set underlying $X^\an$, $\varphi$
induces a bijection of~$|X^\an|$ onto the set $X(\CC)$ of points
of~$X$ with values in $\CC$. Moreover, for each point $x$ of~$X^\an$, the morphism
$$
\varphi_x\colon\cal{O}_{X,\varphi(x)}\to\cal{O}_{X^\an,x},
$$
which is necessarily local, gives, by passing to the
completions, an isomorphism
$$
\widehat{\varphi}_x\colon\widehat{\cal{O}}_{X,\varphi(x)}\isomto
\widehat{\cal{O}}_{X^\an,x},
$$
In particular the morphism $\varphi$ is flat.
\end{theoremdefinition}

Let us note that the fact that $\varphi$ induces a bijection of~$X^\an$ onto~$X(\CC)$ follows from the universal property of~$X^\an$.
On the other hand, we have the following assertions:

\begin{enumerate}
\item [a)] If the theorem is true for a scheme $Y$, the same
holds for every subscheme $X$ of~$Y$. Let us first suppose
that $X$ is an open subscheme of~$Y$; if $\psi\colon
Y^\an\to Y$ is the canonical morphism, $\psi^{-1}(X)$ is an open subset
% original p. 313
of~$Y^\an$ which one endows with the analytic space structure induced
by that of~$Y^\an$. Since every morphism from an analytic space
$\othercal{X}$ into $X$ factors through $Y^\an$ by the
universal property of the latter, hence through
$X^\an$ which is the fiber product $Y^\an\times_Y X$, $X^\an$ is
the analytic space associated with $X$. Finally the assertion
concerning the $\varphi_x$ is evident.

It now suffices to consider the case where $X$ is a
closed subscheme of~$Y$. Let~$I$ be the coherent $\cal{O}_Y$-Ideal
defining $X$; then $I\cdot\cal{O}_{Y^\an}$ is a
coherent sheaf of ideals on~$\cal{O}_{Y^\an}$ which
defines a closed analytic subspace $X^\an$ of~$Y^\an$; one
sees as in the case of an open subscheme that $X^\an$ is
the analytic space associated with $X$. Let $\varphi\colon
X^\an\to X$ be the canonical morphism. For every point $x$ of~$X^\an$, the
morphism $\varphi_x$ is none other than the morphism
$$
\cal{O}_{Y,\psi(x)}/I_{\psi(x)}\to
\cal{O}_{Y^\an,x}/I_{\psi(x)}\cdot\cal{O}_{Y^\an,x}
$$
induced by $\psi_x$; its completion
$$
\widehat{\varphi}_x\colon
\widehat{\cal{O}}_{Y,\psi(x)}/I_{\psi(x)}\cdot\widehat{\cal{O}}_{Y,\psi(x)}
\to\widehat{\cal{O}}_{Y^\an,x}/I_{\psi(x)}\cdot\widehat{\cal{O}}_{Y^\an,x}
$$
is an isomorphism since $\widehat\psi_x$ is one, which
proves a).

\item [b)] If one has two $\CC$-schemes $X_1$, $X_2$, such that
$X_1^\an$ and $X_2^\an$ exist, then the same holds for
$(X_1\times X_2)^\an$. Indeed, let $\varphi_1\colon X_1^\an\to
X_1$, $\varphi_2\colon X_2^\an\to X_2$ be the canonical morphisms, and
$p_1$, $p_2$ the two projections of~$X_1^\an\times X_2^\an$. One
deduces formally from EGA~I~1.8.1 that $X_1\times X_2$ is the
product of~$X_1$ and $X_2$ in the category of locally ringed
spaces; it follows that the morphisms
$\varphi_1\cdot p_1$ and $\varphi_2\cdot p_2$ define a
% original p. 314
morphism $\varphi\colon X_1^\an\times X_2^\an\to X_1\times X_2$ and
that the pair $(X_1^\an\times X_2^\an,\varphi)$ represents the
functor $\othercal{X}\mto\Hom_\CC(\othercal{X},X_1\times X_2)$.

\item [c)] If $\othercal{E}^1$ denotes the affine space of dimension 1,
\ie the topological space $\CC$ endowed with the sheaf of holomorphic
functions, the functor%
{{\renewcommand{\thefootnote}{*}\addtocounter{footnote}{-1}\footnote{\lcrochetbf Added in 2003: $E^1_\CC$ denotes the (algebraic) affine line over $\CC$.\rcrochetbf}}}
$\othercal{X}\mto\Hom_\CC (\othercal{X},E^1_\CC)$
is representable by $\othercal{E}^1$, the canonical morphism
$\varphi\colon\othercal{E}^1\to E^1_\CC$ being the
evident morphism. Indeed, giving a morphism from an analytic space
$\othercal{X}$ into $E^1_\CC$ is equivalent to giving an
element of~$\Gamma(\othercal{X},\cal{O}_\othercal{X})$, which also amounts
to giving a morphism from~$\othercal{X}$ into~$\othercal{E}^1$. One
evidently has a bijection $|\othercal{E}^1|\simeq E^1(\CC)$, and, for
each point $x\in\othercal{E}^1$, the morphism $\widehat\varphi_x$ is
none other than the identity morphism of a ring of formal power series
in one variable over~$\CC$.
\end{enumerate}

One deduces from b) and c) that the theorem is true for
the affine space $E^n_\CC$, $n\geq 0$. Using a), one sees that the same
holds for every affine scheme $X$ locally of finite
type over~$\CC$. If one no longer supposes $X$ affine and if $(X_i)$ is
a covering of~$X$ by affine open subsets, it follows from the
universal property and from a) that the $X_i^\an$ glue together
and thus define the analytic space $X^\an$ associated
with~$X$.

\subsection{}
\label{XII.1.2}
Let $f\colon X\to Y$ be a morphism of~$\CC$-schemes locally of
finite type. If $\varphi\colon X^{\an}\to X$ and $\psi\colon Y^{\an}\to
Y$ are the canonical morphisms, it follows from the
universal property of~$Y^{\an}$ that there exists a unique
morphism $f^{\an}\colon X^{\an}\to Y^{\an}$ such that the diagram
$$
\xymatrix{
X^{\an} \ar[r]\ar[d]_-{f^{\an}}& X\ar[d]^-f\\ Y^{\an} \ar[r]& Y
}
$$
is
% original p. 315
commutative. We have thus defined a functor $\Phi$ from the
category of $\CC$-schemes locally of finite type to the
category of analytic spaces.

The functor $\Phi$ commutes with finite projective limits. Indeed, it
suffices to see that $\Phi$ commutes with fiber products. Now, if
$X$, $Y$, $Z$, are schemes locally of finite type over~$\CC$,
it follows from the fact that $X\times_Z Y$ is the fiber product of
$X$ and $Y$ over~$Z$ in the category of locally ringed
spaces that $X^{\an}\times_{Z^{\an}}Y^{\an}$
satisfies the universal property which characterizes
$(X\times_Z Y)^{\an}$.

\subsection{}
\label{XII.1.3}
Let $X$ be a $\CC$-scheme locally of finite type, $X^{\an}$
the associated analytic space, $\varphi\colon X^{\an}\to X$ the
canonical morphism. If $F$ is an ${\cal{O}}_X$-Module, the inverse
image $\varphi^*F=F^{\an}$ is a sheaf of modules on~${\cal{O}}_{X^{\an}}$. One thus defines a functor from the
category of $\cal{O}_X$-Modules to the category of
Modules on~$X^{\an}$. This functor commutes with inductive limits
(EGA~0~4.3.2). The sheaf $\cal{O}_{X^{\an}}$ being coherent
\cite[\no18~\S 2~th\ptbl 2]{XII.4}, it transforms coherent sheaves into
coherent sheaves (EGA~0~5.3.11). We have moreover:

\begin{subproposition}
\label{XII.1.3.1}
The functor which associates with an $\cal{O}_{X^{\an}}$-Module $F$ its
inverse image $F^{\an}$ on~$X^{\an}$ is exact, faithful,
conservative.
\end{subproposition}

Exactness follows from the fact that the morphism $\varphi\colon
X^{\an}\to X$ is flat \eqref{XIII.1.1}. Let us prove that the functor $F\mto
F^{\an}$ is faithful. Taking exactness into account, it suffices to
show that, if $F^{\an}$ is zero, then so is~$F$. Now,
for every point $x$ of~$X^{\an}$, one then has $F_{\varphi
(x)}\otimes_{\cal{O}_{X,\varphi (x)}}\cal{O}_{X^{\an},x}=0$. The
morphism $\cal{O}_{X,\varphi (x)}\to \cal{O}_{X^{\an},x}$ being
faithfully flat, one has $F_{\varphi (x)}=0$ for every closed
point $\varphi (x)$ of~$X$, and, since $X$ is Jacobson
(EGA~IV~10.4.8), this implies that $F$ is zero.

The
% original p. 316
fact that the functor $F\mto F^{\an}$ is conservative is formal
from exactness and faithfulness.

\section[Comparison of the properties of a scheme]{Comparison of the properties of a scheme and of the associated
analytic \hbox{space}}
\label{XII.2}

\begin{proposition}
\label{XII.2.1}
Let $X$ be a $\CC$-scheme locally of finite type, $X^{\an}$
the associated analytic space, $n$ an integer. Consider the
property $P$ of being
\begin{itemize}
\item[(i)] nonempty
\item[(i')] discrete
\item[(ii)] Cohen-Macaulay
\item[(iii)] $(\mathrm{S}_n)$
\item[(iv)] regular
\item[(v)] $(\mathrm{R}_n)$
\item[(vi)] normal
\item[(vii)] reduced
\item[(viii)] of dimension $n$.
\end{itemize}
Then, for $X$ to have the property $P$, it is necessary and
sufficient that $X^{\an}$ have it.
\end{proposition}

Let $\varphi\colon X^{\an}\to X$ be the canonical
morphism. (i)~follows from the fact that one has $|X^{\an}|=X(\CC )$
\eqref{XIII.1.1} and from the fact that $X$ is Jacobson (EGA~IV~10.4.8). To say that~$X$
(\resp $X^{\an}$) is discrete is equivalent to saying that one has $\dim
X=0$ (\resp $\dim X^{\an}=0$ by \cite[\no19~\S 4~cor\ptbl 6]{XII.4}); (i')
therefore follows from~(viii).

Let $P$ be one of the properties (ii) to (vii). For $X$
to have the property $P$, it is necessary and sufficient that $P$
be satisfied at each closed point of~$X$; indeed, $X$
being excellent (EGA~IV~7.8.6~(iii)), the set of points
% original p. 317
where $X$ satisfies $P$ is open (\loccit) and, if this open set
contains all the closed points, it is equal to the whole of $X$.
To say that $X$ (\resp $X^{\an}$) has the property $P$
is therefore equivalent to saying that, for every point $x$ of~$X^{\an}$,
the local ring $\cal{O}_{X,\varphi (x)}$ (\resp $\cal{O}_{X^{\an},x}$)
has the property $P$. Since the fact that an excellent local ring
has the property $P$ can be seen after passing to the
completion, the proposition follows from the isomorphisms
$\widehat{\cal{O}}_{X,\varphi(x)}\isomto\widehat{\cal{O}}_{X^{\an},x}$
in cases (ii) to (vii). The same holds in case
(viii), taking into account the relations
$$
\dim X=\sup_x~\dim \cal{O}_{X,\varphi (x)}\qquad \dim
X^{\an}=\sup_x~\dim \cal{O}_{X^{\an},x}
$$
where $x\in X^{\an}$. This completes the proof.

\begin{proposition}
\label{XII.2.2}
Let $X$ be a $\CC$-scheme locally of finite type,
$\varphi\colon X^{\an}\to X$ the canonical morphism, $T$ a
locally constructible subset of~$X$. Then one has the relation
$$
\varphi^{-1}(\overline{T})=\overline{\varphi^{-1}(T)}.
$$
\end{proposition}

One may suppose that $T$ is a dense open subset of~$X$. Let $H$ be the
reduced closed subscheme of~$X$ with underlying space
$X-T$; the associated space $H^{\an}$ is a closed analytic
subspace of~$X^{\an}$ with underlying space
$X^{\an}-\varphi^{-1}(T)$. One must show that every point $x$ of
$H^{\an}$ belongs to $\overline{\varphi^{-1}(T)}$. Now, at such a
point $x$, the germ of analytic space $(X^{\an},x)$ contains the
subgerm $(H^{\an},x)$, and the latter is defined by a
non-nilpotent Ideal of~$\cal{O}_{X^{\an},x}$. It then follows from the
Nullstellensatz \cite[\no19~\S 4~cor\ptbl 3]{XII.4} that every open neighborhood of~$x$
contains points of~$X^{\an}$ which do not belong to
$H^{\an}$, which indeed proves that one has
$x\in\overline{\varphi^{-1}(T)}$.

\begin{corollary}
\label{XII.2.3}
Let
% original p. 318
$X$ be a $\CC$-scheme locally of finite type, $\varphi
\colon X^\an \to X$ the canonical morphism, $T$ a locally
constructible subset of~$X$. For $T$ to be an open subset (\resp a
closed subset, \resp a dense subset), it is necessary and sufficient that
$\varphi^{-1} (T)$ be so.
\end{corollary}

The corollary follows from~\Ref{XII.2.2} and from the fact that, $X$
being a Jacobson scheme (EGA~IV~10.4.8), two
locally constructible subsets of~$X$ which have the same trace on the
very dense set $X (\CC )$ are equal.

\begin{proposition}
\label{XII.2.4} Let $X$ be a $\CC$-scheme locally of
finite type. For $X$ to be connected (\resp irreducible), it is necessary
and sufficient that $X^\an$ be so.
\end{proposition}

Suppose $X^\an $ connected (\resp irreducible). The image $X (\CC)
$ of~$X^\an$ in $X$ is then connected (\resp irreducible). It
follows that $X$ is connected (\resp irreducible) since the
closed subsets of~$X$ and of $X (\CC)$ correspond bijectively
(EGA~IV~10.1.2).

Conversely, suppose $X$ connected (\resp irreducible), and
let us show that the same holds for~$X^\an$. One may restrict oneself to the case
where $X$ is irreducible. Indeed, suppose $X$ connected.
Given a point $x$ of~$X$, the set of points $y \in X$
such that there exists a finite sequence of irreducible closed
subschemes $X_1, \dots, X_n$ of~$X$, with $x \in X_1$, $y \in
X_n$, $X_i \cap X_{i+1} \neq \emptyset $ for $1 \leq i \leq n-1$,
is a set which is both open and closed, hence equal to
the whole of $X$. For a sequence $X_1, \dots, X_n$ as
above, one also has $X_{i}^\an \cap X_{i+1}^\an \neq
\emptyset $ for $1 \leq i \leq n-1 $; if one supposes it
proved that the $X_{i}^\an$ are connected, then the same
holds for~$X^\an$.

From
% original p. 319
now on we suppose $X$ irreducible. One may moreover suppose
$X$ affine. Indeed, if $(U_i)_{i \in I}$ is a covering of~$X$
by affine open subsets, any two of these open sets have a nonempty
intersection, and the same property is therefore true for the
covering $(U_{i}^\an)_{i \in I} $ of~$X^\an$; if one supposes it
proved that the $U_{i}^\an $ are irreducible, then the same
holds for~$X^\an$.

One may moreover suppose that $X$ is normal. Indeed, let $\tilde X$ be
the normalization of~$X$; since the morphism $\tilde X \to X$ is
surjective, so is $\tilde X^\an \to X^\an$, which
proves that, if $\tilde X^\an $ is irreducible, so
is~$X^\an$.

From now on we suppose $X$ affine and normal. Since the local rings
of~$X^\an$ are integral, it amounts to the same to say that
$X^\an$ is irreducible or that it is connected. Indeed, if
$\cal{F}$ is a closed analytic subset of~$X^\an$, the set
of points $x$ of~$X^\an$ where one has $\codim_x (\cal{F}, X^\an ) =
0$ is a closed analytic subset of~$X^\an$ \cite[\no 20 A
Cor\ptbl 1]{XII.4} which is also open; if $X^\an$ is connected, this proves that,
if one has $\cal{F} \neq X^\an$, $\cal{F}$ is nowhere dense, hence that $X^\an$
is irreducible. One is thus reduced to showing that
$X^\an$ is connected.

Let
$$
i \colon X \to P
$$
be a compactification of~$X$, where $P$ is a normal projective
$\CC$-scheme and $i$ a dominant open immersion. It then
follows from \cite[\no 12 th\ptbl 1]{XII.10} that $P^\an$ is connected. Since
$X^\an$ is obtained by removing from $P^\an$ a nowhere dense closed
analytic subset, it follows from~\Ref{XII.2.5} below that $X^\an$
is also connected.

\begin{lemma}
\label{XII.2.5}
Let
% original p. 320
$\cal{P} $ be a connected normal analytic space, $\cal{Y}$ a
nowhere dense closed analytic subset; then $\cal{X} = \cal{P} - \cal{Y}$
is connected.
\end{lemma}

When $\cal{Y}$ is of codimension $\geq 2$, the proposition
follows from \cite[\no 3 prop\ptbl 4]{XII.11}. In the general case one may
suppose, at the cost of removing from $\cal{P}$ a closed analytic
subset of codimension $\geq 2$, that $\cal{P}$ and $\cal{Y}$
(considered as a reduced analytic subspace of~$\cal{P}$) are regular.
By the implicit function theorem, every point $y$ of~$\cal{Y}$ has a
neighborhood $\cal{U}$ isomorphic to a ball of an affine space
$\cal{E}^{n}$, such that $\cal{U} \cap \cal{Y}$ is defined by
the vanishing of a certain number of coordinate functions. This
proves that $\cal{U} - \cal{U} \cap \cal{Y}$ is connected, and
therefore so is~$\cal{X}$.

\begin{corollary}
\label{XII.2.6} Let $X$ be a $\CC$-scheme locally of finite
type; the morphism
$$
\pi_0 (X^{\an}) \to \pi_0 (X)
$$
induced by the canonical morphism $X^\an \to X$ is bijective.
\end{corollary}
